% Focused independent analytic audit, 2026-10-08. Old files unchanged.
\subsection{The individual rough-factor Siegel--Walfisz bridge}
\label{sec:focused-rough-sw-product-interface-audit}

The chosen interface is the additional roughness condition on the single
Heath--Brown variable used in L23 Proposition 5.2. L23 Definition 2.1
requires all auxiliary integers, not merely a Siegel--Walfisz estimate
with no auxiliary coprimality. Its examples $1$ and $\mu$ do not by
themselves cover the rough truncation. This note independently derives
that bridge from classical inputs, and identifies the precise limits of
the conclusion. It is not a verification of all spectral arguments or
of the complete new-$K$ theorem.

\paragraph{Source status and what is actually being quoted.}
Here L23 is the versioned primary source
\href{https://arxiv.org/abs/2309.08522v1}{arXiv:2309.08522v1}.
The local extraction \texttt{external-source-audit/lichtman2023v1.txt}
has Definition 2.1 at lines 388--396 and Proposition 5.2 at
946--962. Its requirement is
\[
 \left|\sum_{\substack{n\sim U\,,\ n\equiv r\ (q)\\(n,D)=1}}
 \alpha_n-\frac1{\varphi(q)}\sum_{\substack{n\sim U\\(n,Dq)=1}}
 \alpha_n\right|
 \ll_A U\tau(D)^{O(1)}(\log U)^{-A}
\]
for every $D,q\ge1$ and $(r,q)=1$. Its convolution proposition
explicitly takes a supremum over $0<|a_{\rm res}|<x^{1+\epsilon}$.
Thus the source is not restricted to one residue fixed independently
of $x$.

The arXiv record checked on 8 October 2026 exposes v1 and no journal
reference. This audit has not verified journal publication of L23 and
does not describe it as a checked journal theorem. This bibliographic
qualification does not itself refute its mathematics. Maynard's
\emph{Primes in Arithmetic Progressions to Large Moduli II} is listed as
a 2025 publication by his
\href{https://www.maths.ox.ac.uk/research/faculty-books/authors-k-o/james-maynard}
{institutional author bibliography}; the local argument's theorem
numbers are from \href{https://arxiv.org/abs/2006.07088v1}{its v1 source}.
This audit does not silently identify its page or proposition numbers
with those of the published edition.

The published classical inputs used in the proof below are the ordinary
M\"obius Siegel--Walfisz estimate, the dimension-one fundamental lemma,
and fixed-order divisor moments. L23 itself records the unrough
M\"obius example and cites Friedlander--Iwaniec,
\emph{Opera de Cribro}, Theorem 6.12, for its fundamental-lemma use.
The rough extension proved next is a derived lemma of this project,
not the literal statement of an external theorem.

\paragraph{Statement with all relevant quantifiers.}
Fix $\xi>0$ and $B\ge0$. Let
\[
 x^\xi\le U\le x,\qquad Z=x^{1/(\log\log x)^3},
 \qquad I\subset[U,2U].
\]
For $e=1$ or $e=\mu$, set
\[
 a(n)=e(n)\mathbf1_{P^-(n)>Z}\mathbf1_I(n)h(n)n^{it},
\]
where $|h|\le1$, $\int_U^{2U}|h'(v)|\,dv\le C_h(\log x)^B$,
and $|t|\le(\log x)^B$ with fixed $C_h$.
Then for every $A>0$, every integer $D,q\ge1$, and $(r,q)=1$,
\begin{equation}\label{eq:focused-rough-sw-conclusion}
 \left|\sum_{\substack{n\equiv r\ (q)\\(n,D)=1}}a(n)
 -\frac1{\varphi(q)}\sum_{(n,Dq)=1}a(n)\right|
 \ll_{A,B,\xi,C_h} U\tau(D)(\log x)^{-A}.
\end{equation}
Intervals and weights may vary with $x$ or with the even target $N$,
provided these explicit bounds are uniform. There is no size restriction
on $D$, no exceptional set of target integers, and no assumption of a
Siegel--Walfisz theorem for a grouped convolution.

\paragraph{Step 1: auxiliary coprimality for the unrough M\"obius sum.}
For any $E\ge1$, the coefficient identity
\[
 \mu(n)\mathbf1_{(n,E)=1}
 =\sum_{\substack{v\mid n\\v\mid E^\infty}}\mu(n/v)
\]
holds prime by prime. On an interval in $[V,2V]$, split at
$v=V^{1/4}$. For the smaller $v$, ordinary M\"obius
Siegel--Walfisz applies on a scale at least $V^{3/4}$, uniformly for
any fixed logarithmic range of $q$. Terms $(v,q)>1$ vanish in both
primitive-residue expressions. Moreover
\[
 \sum_{v\mid E^\infty}\frac1v
 =\prod_{p\mid E}(1-1/p)^{-1}\le\tau(E).
\]
The absolute mass of the tail, including the normalized coprime mean,
is bounded by
\[
 O\left(V\sum_{\substack{v>V^{1/4}\\v\mid E^\infty}}\frac1v\right)
 \ll V^{7/8}\prod_{p\mid E}(1-p^{-1/2})^{-1}
 \le V^{7/8}\tau(E)^2.
\]
The last product inequality holds factor by factor, including $p=2$.
Thus for any fixed $C,B_0$ the unrough auxiliary discrepancy obeys
\begin{equation}\label{eq:focused-unrough-auxiliary}
 \mathcal E_E(V;q,r)
 \ll_{C,B_0} V\tau(E)(\log V)^{-C}
                   +V^{7/8}\tau(E)^2
\end{equation}
uniformly for $q\le(\log V)^{B_0}$ and for every interval endpoint.
The divisor exponent is fixed; it does not grow with $E$.

\paragraph{Step 2: impose roughness, including the signed error.}
First take $D=1$, $h=1$, $t=0$, and $q\le(\log x)^{B_0}$.
Use upper and lower fundamental-lemma polynomials
$\rho^\pm(n)=\sum_{\ell\mid n}\lambda_\ell^\pm$ of level
$L=U^{1/4}$, with $|\lambda_\ell^\pm|\le1$ and squarefree
$\ell\mid P(Z)$. On a primitive class the primes dividing $q$ have
local density zero. Deleting them preserves the dimension-one bound;
one can equivalently choose the auxiliary polynomials on the deleted
prime set for this proof. Their constants are uniform in $q$.
With $s_0=\log L/\log Z\ge(\xi/4)(\log\log x)^3$, elementary
counting of multiples and the fundamental lemma give
\[
 \sum_{\substack{n\in I\\n\equiv r\ (q)}}(\rho^+-\rho^-)(n)
 \ll (U/q)e^{-c s_0}+L,
\]
and the same bound for the average of primitive classes. This is
$O_C(U(\log x)^{-C})$ for every fixed $C$.

The signed case is justified by the absolute inequality
\[
 |\mu(n)\mathbf1_{P^-(n)>Z}-\mu(n)\rho^+(n)|
 \le\rho^+(n)-\rho^-(n)
\]
on those classes. It does not compare the signed sums by positivity.
For a squarefree $\ell$,
$\mu(\ell m)=\mu(\ell)\mu(m)\mathbf1_{(m,\ell)=1}$.
After expanding the polynomial, apply
\eqref{eq:focused-unrough-auxiliary} with $E=\ell$ and $V=U/\ell$.
Here $\ell\le U^{1/4}$ and $V\ge U^{3/4}$, so the same logarithmic
modulus range applies to every inner sum. The two aggregate errors are
\begin{align*}
 &\ll_C \frac{U}{(\log x)^C}
          \sum_{\ell\le L}\frac{\tau(\ell)}\ell
       \ll_C U(\log x)^{2-C},\\
 &\ll U^{7/8}\sum_{\ell\le L}\frac{\tau(\ell)^2}{\ell^{7/8}}
       \ll U^{29/32}(\log U)^4.
\end{align*}
Both permit any prescribed logarithmic saving, since $U\ge x^\xi$.
This explicit power estimate closes the potential hidden auxiliary
coprimality loss in the signed rough expansion. For $e=1$, each
expanded progression count differs from its primitive-class average
by $O(1)$, so the corresponding error is merely $O(L)$.

\paragraph{Step 3: all moduli and all auxiliary integers.}
For a one-bounded sequence and $D=1$, the trivial discrepancy is
\[
 O(U/q+1+U/\varphi(q))\ll Uq^{-1/2}+1.
\]
Taking $B_0$ sufficiently large handles every
$q>(\log x)^{B_0}$ with the requested saving.
Now an arbitrary $D$ can delete only rough multiples of primes
$p\mid D$ with $p>Z$. Their total absolute number in $[U,2U]$ is
at most $(U/Z+1)\omega(D)$. Thus its effect on the discrepancy is at
most twice this quantity, and
\[
 (U/Z+1)\omega(D)\ll_{A,\xi}U\tau(D)(\log x)^{-A}.
\]
This is uniform even when $D$ is much larger than every other scale.
Finally $h(v)v^{it}$ has total variation
$O((\log x)^B)$, so partial summation, beginning with a stronger
logarithmic saving, proves \eqref{eq:focused-rough-sw-conclusion}.

\paragraph{Two adversarial boundary tests.}
The $\tau(D)$ allowance cannot be silently removed. For $e=1$, let
$I=(U,2U]$ and let $D$ be the product of all primes in $I$ congruent
to one modulo three. These primes exceed $Z$; in $I$, imposing
$(n,D)=1$ deletes exactly those primes. If their number is $P_U$, the
discrepancy for $q=3,r=1$ changes by $-P_U/2$. The base discrepancy
is smaller than any inverse logarithmic power, whereas classical PNT
in that progression gives $P_U\sim U/(2\log U)$. A hypothetical bound
$O_A(U(\log x)^{-A})$ independent of $D$ would therefore fail.
The actual lemma has the required $\tau(D)$ and is not contradicted.

Nor can a unit Heath--Brown slot be selected as the large variable.
For the rough sequence supported just at $n=1$, take $q=3$, $r=1$,
$D=1$; its discrepancy is $1/2$, not $O_A((\log x)^{-A})$.
The actual Type II selection uses $U\ge x^{2\zeta}$ and excludes this
case. No assertion uniform as $\xi\downarrow0$ has been proved.

\paragraph{Map back to the new modulus tree.}
The admissible product family fixes a least outer exponent $a>0$ and
a finite number of original prime slots. Choose $12\zeta<a/2$ after
fixing buffered support margins. Apply the $k=3$ Heath--Brown identity
at each original slot scale, and retain that original scale before
Mellin separation. A selected individual medium variable satisfies
\[
 x^{2\zeta}\le U\le x^{1/3+\zeta/2},
\]
and its coefficient is precisely rough $1$ or rough $\mu$, times an
interval, a normalized logarithm, and a bounded-variation Mellin twist.
Set $\xi=2\zeta$ in the proved lemma. This verifies every part of
L23 Definition 2.1 for the selected sequence, uniformly in the local
scale and target $N$. The complementary convolution is not required
to satisfy this property; its separate divisor truncation remains a
different step.

The inspected fourth-prefix rule freezes a new level at even rank four
and changes later odd tests only. The fifth rule changes the odd fifth
decision and freezes a level afterwards. The sixth rule keeps rank-six
membership and changes later odd tests only, using an upper bound for
the actual full fifth cap. All use positive fixed cutoffs, finite
partitions, and full plus selected-remaining square gates. These
choices modify the modulus coefficients, not the individual outer
Heath--Brown variable. They cannot add a new arithmetic restriction to
that variable. For each finite modulus atom the group coefficients may
depend on $N$, but L23 permits arbitrary one-bounded complex group
coefficients with the stated coprimality. Such dependence is harmless
once the atom's partition and factorization inequalities are justified.

For the target residue, local product scales
$x\ge N^{1-\kappa}/O(1)$ with
$\kappa<\epsilon/(2(1+\epsilon))$ give $N<x^{1+\epsilon}$ for large
$N$. L23's displayed supremum then includes every target $N$, with no
averaging over target integers. This observation and the proved SW
lemma remove no modulus-factor size inequality: those remain required
for every atom. Shrinking the fixed outer cutoff changes $\zeta$ and
implicit constants, not the limiting tree density or its $K$.

\paragraph{Precisely what remains outside this focused result.}
The proof above closes the rough-SW bridge; no additional rough-SW
hypothesis is needed for any fixed positive outer cutoff. It does not
turn the numerical inequality $K_{\rm end}<6.7619$ into a quoted
external theorem. The following are separate project-derived
propositions, with local proof fragments rather than published theorem
statements: the nonleading selected-subproduct factorization and its
capacity repair; compilation of the fourth/fifth/sixth endpoint strategy
into finite signed box atoms; the uniform upper-density evaluation of
that tree; and the positive radial majorant for sharp product domains.
Their existence is not established by a finite rational box sum alone.
This audit inspected their role and support changes but did not reprove
all of them or the external dispersion theorem.

Accordingly the concrete conclusion is narrower than a global PASS:
the full auxiliary-coprimality and twist hypothesis of L23 for the
individual rough Type II factor is now supplied by the explicit proof
above. No counterexample to that hypothesis, or new restriction on the
fixed positive-cutoff class, was found. The every-even-$N$ new-$K$
conclusion still requires the separate signed-atom/factorization,
density, exceptional one/two-long-variable, and radial-limit arguments.
No claim for arbitrary arithmetic outer weights, vanishing cutoffs,
or infinitely adaptive prefix strategies is licensed.
